% basic_principles_elucidations.tex -- Book I. Elucidations (jemmybutton-matched).
% Text and figures from jemmybutton byrne-en-latex.tex (CC-BY-SA).
% Build: lualatex basic_principles_elucidations.tex  ->  basic_principles_elucidations.pdf
\documentclass[12pt]{article}
\usepackage{geometry}
\geometry{a5paper, margin=1.5cm}
\usepackage[lining]{ebgaramond}
\usepackage{amssymb}
\usepackage{byrne}
\def\mpPre{u := 1cm; textLabels := false;}

\begin{document}

\begin{center}
  {\large\textbf{Book~I.}}\\[4pt]
  {\Large\textbf{Elucidations}}
\end{center}

\vspace{0.4cm}
\noindent\rule{\linewidth}{0.4pt}
\vspace{0.3cm}

\noindent
Geometry has for its principal objects the exposition and explanation of the
properties of \emph{figure}, and figure is defined to be the relation which
subsists between the boundaries of space. Space or magnitude is of three kinds,
\emph{linear}, \emph{superficial}, and \emph{solid}.

\vspace{0.5cm}

% --- Figure 1: simple angle at vertex A ---
\defineNewPicture{%
  pair A, B, C;%
  numeric s;%
  s := 3/2u;%
  A := (0, s);%
  B := (1/2s, 0);%
  C := B xscaled -1;%
  draw byAngle.A(B, A, C, byyellow, SOLID_SECTOR);%
  byLineDefine(B, A, byblue, SOLID_LINE, REGULAR_WIDTH);%
  byLineDefine(C, A, byred,  SOLID_LINE, REGULAR_WIDTH);%
  draw byNamedLineSeq(0)(CA,BA);%
  label.urt(btex A etex, A);%
}
\noindent
\begin{minipage}[t]{0.28\linewidth}
  \vspace{0pt}\centering
  \drawCurrentPicture
\end{minipage}\hfill
\begin{minipage}[t]{0.68\linewidth}
  \vspace{0pt}
  Angles might properly be considered as a fourth species of magnitude.
  Angular magnitude evidently consists of parts, and must therefore be
  admitted to be a species of quantity. The student must not suppose that
  the magnitude of an angle is affected by the length of the straight lines
  which include it and of whose mutual divergence it is the measure.
  The \emph{vertex} of an angle is the point where the \emph{sides} or
  \emph{legs} of the angle meet, as~A.
\end{minipage}

\vspace{0.6cm}

% --- Figure 2: multiple angles at vertex C ---
\defineNewPicture{%
  pair B, C, D, E, F, G, H;%
  numeric s;%
  s := 5/4u;%
  C := (0, 0);%
  B := dir(0)*s;%
  D := dir(50)*s;%
  E := dir(-30)*s;%
  F := E scaled -1;%
  G := D scaled -1;%
  H := B scaled -1;%
  angleScale := 4/3;%
  draw byAngle(E, C, B, byyellow, SOLID_SECTOR);%
  draw byAngle(B, C, D, byblack,  SOLID_SECTOR);%
  draw byAngle(D, C, F, byblue,   SOLID_SECTOR);%
  draw byAngle(F, C, H, byred,    SOLID_SECTOR);%
  draw byAngle(H, C, G, byyellow, ARC_SECTOR);%
  draw byAngle(G, C, E, byblue,   ARC_SECTOR);%
  draw byLine(B, H, byblue,  SOLID_LINE, REGULAR_WIDTH);%
  draw byLine(D, G, byred,   SOLID_LINE, REGULAR_WIDTH);%
  draw byLine(E, F, byblack, SOLID_LINE, REGULAR_WIDTH);%
  label.bot(btex C etex, C shifted (0, -3pt));%
  label.bot(btex B etex, B);%
  label.lrt(btex D etex, D);%
  label.llft(btex F etex, F);%
  label.bot(btex H etex, H);%
  label.lrt(btex G etex, G);%
  label.llft(btex E etex, E);%
}
\noindent
\begin{minipage}[t]{0.34\linewidth}
  \vspace{0pt}\centering
  \drawCurrentPicture
\end{minipage}\hfill
\begin{minipage}[t]{0.62\linewidth}
  \vspace{0pt}
  An angle is often designated by a single letter when its legs are the only
  lines which meet together at its vertex. Thus the red and blue lines form
  the yellow angle, which in other systems would be called angle~A. But when
  more than two lines meet in the same point, it was necessary by former
  methods, in order to avoid confusion, to employ three letters to designate
  an angle about that point, the letter which marked the vertex of the angle
  being always placed in the middle. Thus the black and red lines meeting
  together at~C, form the blue angle, and has been usually denominated the
  angle~FCD or~DCF. The lines~FC and~CD are the legs of the angle; the
  point~C is its vertex. In like manner the black angle would be designated
  the angle~DCB or~BCD. The red and blue angles added together, or the
  angle~HCF added to~FCD, make the angle~HCD; and so of other angles.
\end{minipage}

\vspace{0.5cm}

\noindent
When the legs of an angle are produced or prolonged beyond its vertex, the
angles made by them on both sides of the vertex are said to be
\emph{vertically opposite} to each other: thus the red and yellow angles are
said to be vertically opposite angles.

\vspace{0.4cm}

\noindent
\emph{Superposition} is the process by which one magnitude may be conceived
to be placed upon another, so as exactly to cover it, or so that every part
of each shall exactly coincide.

\vspace{0.4cm}

\noindent
A line is said to be \emph{produced}, when it is extended, prolonged, or it
has its length increased, and the increase of length which it receives is
called the \emph{produced part}, or its \emph{production}.

\vspace{0.4cm}

\noindent
The entire length of the line or lines which enclose a figure, is called its
\emph{perimeter}. The first six books of Euclid treat of plane figures only.
A line drawn from the centre of a circle to its circumference, is called a
\emph{radius}. That side of a right angled triangle, which is opposite to the
right angle, is called the \emph{hypotenuse}. An oblong is defined in the
second book, and called a \emph{rectangle}. All lines which are considered in
the first six books of the Elements are supposed to be in the same plane.

\vspace{0.4cm}

\noindent
The \emph{straight-edge} and \emph{compasses} are the only instruments, the
use of which is permitted in Euclid, or plain Geometry. To declare this
restriction is the object of the \emph{postulates}.

\vspace{0.4cm}

\noindent
The \emph{Axioms} of geometry are certain general propositions, the truth of
which is taken to be self-evident and incapable of being established by
demonstration.

\vspace{0.4cm}

\noindent
\emph{Propositions} are those results which are obtained in geometry by a
process of reasoning. There are two species of propositions in geometry,
\emph{problems} and \emph{theorems}.

\vspace{0.4cm}

\noindent
A \emph{Problem} is a proposition in which something is proposed to be done;
as a line to be drawn under some given conditions, a circle to be described,
some figure to be constructed, \&c. The \emph{solution} of the problem
consists in showing how the thing required may be done by the aid of the
straight-edge and compasses. The \emph{demonstration} consists in proving
that the process indicated in the solution attains the required end.

\vspace{0.4cm}

\noindent
A \emph{Theorem} is a proposition in which the truth of some principle is
asserted. This principle must be deduced from the axioms and definitions, or
other truths previously and independently established. To show this is the
object of demonstration.

\vspace{0.5cm}

\begin{center}
A \emph{Problem} is analogous to a postulate.\\[4pt]
A \emph{Theorem} resembles an axiom.\\[4pt]
A \emph{Postulate} is a problem, the solution to which is assumed.\\[4pt]
An \emph{Axiom} is a theorem, the truth of which is granted without demonstration.
\end{center}

\vspace{0.5cm}

\noindent
A \emph{Corollary} is an inference deduced immediately from a proposition.

\vspace{0.4cm}

\noindent
A \emph{Scholium} is a note or observation on a proposition not containing an
inference of sufficient importance to entitle it to the name of
\emph{corollary}.

\vspace{0.4cm}

\noindent
A \emph{Lemma} is a proposition merely introduced for the purpose of
establishing some more important proposition.

\end{document}
